% concatenated from Main.tex
\documentclass[10pt, A4paper]{amsart}
%\makeindex
\usepackage[hidelinks]{hyperref}
% Neu hinzugefügt
\usepackage{cleveref}

\usepackage{amssymb,amsmath,mathtools}
\usepackage{enumerate}
\usepackage{color}
\usepackage{mathrsfs}
\usepackage{enumitem}
\usepackage{graphicx}
\usepackage{float}
\usepackage{tikz}
\usepackage{mathtools}

%\usepackage[notref, notcite]{showkeys}%%%%#####




\usepackage{setspace}

\usepackage[a4paper, left=25mm, right=25mm, top=25mm, bottom=25mm]{geometry}


%\linespread{1.0}

%\usepackage{environ}
%\NewEnviron{killcontents}{}
%\let\proof\killcontents
%\let\endproof\endkillcontents

\newtheorem{theorem}{Theorem}[section]
\newtheorem{proposition}[theorem]{Proposition}
\newtheorem{lemma}[theorem]{Lemma}
\newtheorem{corollary}[theorem]{Corollary}
\theoremstyle{definition}
\newtheorem{definition}[theorem]{Definition}
\newtheorem{example}[theorem]{Example}
\newtheorem{remark}[theorem]{Remark}
\newtheorem{remarks}[theorem]{Remarks}
\newtheorem{assumption}[theorem]{Assumption}

\newcommand{\sint}{\scalebox{1.3}{$\int$}}



\newcommand{\indep}{\perp \!\!\! \perp}
%\newcommand{\marriage}{\circ \!\!\! \circ}
\newcommand{\marriage}{\tikz[baseline=-0.25em, line width=.04em]{%
\draw (0,0) circle[radius=.20em];
\draw (0.20em,0) circle[radius=.20em];
}}


\DeclareFontFamily{U}{mathx}{}
\DeclareFontShape{U}{mathx}{m}{n}{ <-> mathx10 }{}
\DeclareSymbolFont{mathx}{U}{mathx}{m}{n}
\DeclareFontSubstitution{U}{mathx}{m}{n}

\DeclareMathAccent{\widecheck}{0}{mathx}{"71}

\newcommand{\MB}[1]{{\marginpar{\color{purple}{[MB: #1]}}}}
\newcommand{\MZ}[1]{{\marginpar{\color{blue}{[MZ: #1]}}}}
\newcommand{\MBlong}[1]{{{\small\color{purple}{[MB: #1]}}}}
\newcommand{\MZlong}[1]{{{\small\color{blue}{[MZ: #1]}}}}


\makeatletter
\newcommand{\wcheck}[1]{\mathpalette\wcheck@{#1}}
\newcommand{\wcheck@}[2]{%
\begingroup
\edef\wcheck@font{\the
\ifx#1\displaystyle\textfont\else\ifx#1\textstyle\textfont
\else\ifx#1\scriptstyle\scriptfont\else\scriptscriptfont\fi\fi\fi\@ne
}%
\sbox\z@{\wcheck@font\mbox{#2}\mbox{\char\the\skewchar\font}}%
\sbox\tw@{\wcheck@font#2\char\the\skewchar\font}%
\dimen@=\dimexpr\wd\tw@-\wd\z@\relax
{\,\kern2\dimen@\widecheck{\!\kern-2\dimen@#2\!}\,}%
\endgroup
}
\makeatother


%\newcommand{\comment}[1]{\marginpar{ {\color{black}#1} } }
%\newcommand{\comment}[1]{ \marginpar{{#1}}{ \color{black}{\Huge $\bullet$}} }
%\newcommand{\comment}[1]{\marginpar{\color{black}#1}{\color{black}{\Huge $\bullet$}} }
%\newcommand{\comment}[1]{\marginpar{\color{black}#1}{\color{red}$(\ast\ast)$} }
\newcommand{\comment}[1]{}
\newcommand{\F}{\mathcal F}
\newcommand{\G}{\mathcal G}
\newcommand{\cH}{\mathcal H}
\newcommand{\I}{\mathcal I}
\newcommand{\cHs}{\mathcal H}
\newcommand{\AW}{\mathcal A\mathcal W}
\newcommand{\Law}{\mathscr L}
\newcommand{\law}{\Law}

\newcommand{\CFP}{\mathrm{CFP}}
\newcommand{\CFS}{\mathrm{CFS}}
\newcommand{\ip}{\mathrm{ip}}
\newcommand{\id}{\mathrm{id}}
\newcommand{\pp}{\mathrm{pp}}
\newcommand{\pj}{\text{pj}}
\newcommand{\simhk}{\sim_{\mathrm{HK}}}

\newcommand{\cS}{\mathcal{S}}

%\newcommand{\fp}[1]{{\bf #1}}
%\newcommand{\cfp}[1]{{\overline{\bf #1}}}

\newcommand{\fp}[1]{{\mathbb #1}}
\newcommand{\cfp}[1]{{\overline{\mathbb #1}}}
\newcommand{\sawu}[1]{{\widecheck{#1}}}
\newcommand{\fpsawu}[1]{{\widecheck{\mathbb #1}}}
\newcommand{\iuc}[1]{{#1^\sharp}}
\newcommand{\sr}[1]{{\underline{#1}}} %skorokhod representative

\newcommand{\W}{\mathcal W}
\newcommand{\N}{\mathbb N}
\newcommand{\R}{\mathbb R}
\newcommand{\X}{\mathcal X}
\newcommand{\Y}{\mathcal Y}
\newcommand{\Z}{\mathcal Z}
\newcommand{\A}{\mathcal A}
\newcommand{\B}{\mathcal B}
\newcommand{\Pc}{\mathcal P}
\newcommand{\proj}{\mathrm{pj}}
\newcommand{\aproj}{\mathrm{apj}}
\newcommand{\Geod}{\mathrm{Geod}}
\newcommand{\sip}[1]{\overline{\mathcal#1}}

\newcommand{\cpl}{\mathrm{cpl}}
\newcommand{\cplc}{\cpl_{\mathrm{c}}}
\newcommand{\cplbc}{\cpl_{\mathrm{bc}}}

\newcommand{\scpl}{\mathrm{cpl}}
\newcommand{\scplc}{\scpl_{\mathrm{c}}}
\newcommand{\scplbc}{\scpl_{\mathrm{bc}}}

\newcommand{\fpcpl}{\mathrm{PrCpl}}

\newcommand{\dd}{\mathrm{d}}




\newcommand{\unfold}{\mathrm{uf}}
\newcommand{\AF}{\mathrm{AF}}
\newcommand{\depth}{\mathrm{depth}}
\newcommand{\CW}{\mathcal C\mathcal W}
\newcommand{\SCW}{\mathcal S\mathcal C\mathcal W}
\newcommand{\SAW}{\mathcal{SAW}}

\newcommand{\CD}{\mathcal{CW}}
\newcommand{\WCD}{\mathcal{CW}_{w}}
\newcommand{\WAD}{\mathcal{AW}_{w}}
\newcommand{\SWCD}{\mathcal{SCW}_{w}}
\newcommand{\SWAD}{\mathcal{SAW}_{w}}

\def\P{{\mathbb P}}
\def\E{{\mathbb E}}

\def\PP{{\mathcal P}}
\def\Q{{\mathbb Q}}
\def\R{{\mathbb R}}
\newcommand{\Extw}{\mathcal E}
%\newcommand{\FP}{\mathscr{FP}}
\newcommand{\FP}{\mathcal{FP}}
\newcommand{\FFP}{\mathrm{FP}}
\newcommand{\WA}{\mathcal W}


\newcommand{\OB}{Ob{\l}{\'o}j}%{Ob{\l}{\'o}j} 
\newcommand{\HL}{Henry-Labord{\`e}re}


\usepackage{relsize}
%\usepackage{oz}%\fcmp
\def\fcmp{\mathbin{\raise 0.6ex\Hbox{\oalign{\cHfil$\scriptscriptstyle \mathrm{o}$\Hfil\cr\Hfil$\scriptscriptstyle\mathrm{9}$\Hfil}}}}
\let\mcmp\fcmp


%\marginpar{#1}{ aaa }
\numberwithin{equation}{section}

\author{Mathias Beiglb\"ock} 
\author{Markus Zona}
%\thanks{The authors acknowledge FWF-support through project P35197.}
%\author{you}
\keywords{causal transport, adapted Wasserstein distance}
\date{}
%\subjclass[2010]{XXX}



%---------------------------------------------------------------------------------

\begin{document}
\title[Adapted Pinsker]{Pinsker's inequality for adapted total variation}



\begin{abstract} Pinsker's classical inequality asserts that the total variation $TV(\mu, \nu)$ between two probability measures is bounded by $\sqrt{ 2H(\mu|\nu)}$ where $H$ denotes the relative entropy (or Kullback-Leibler divergence). 
Considering the discrete metric, $TV$ can be seen as a Wasserstein distance and as such possesses an \emph{adapted} variant $ATV$. Adapted Wasserstein distances have distinct advantages over their classical counterparts when $\mu, \nu$ are the laws of stochastic processes $(X_k)_{k=1}^n, (Y_k)_{k=1}^n$ and exhibit numerous applications from stochastic control to machine learning. In this note we observe that the adapted total variation distance $ATV$ satisfies the Pinsker-type inequality
$$ ATV(\mu, \nu)\leq \sqrt{n} \sqrt{2 H(\mu|\nu)}.$$
\end{abstract}



\keywords{causal optimal transport, adapted Wasserstein distance, Pinsker's inequality}


\maketitle

\section{Introduction}%
Throughout $\mu, \nu$ denote probabilities on some Polish space $\X$. Their total variation distance 
is 
\begin{align*}
TV(\mu,\nu)= %\sup_{|f|\leq 1} \int f\, \dd (\mu-\nu) =
%|\mu\vee \nu|-|\mu\wedge\nu|.
\int{\textstyle  \left| \frac{\dd \mu}{\dd(\mu+\nu)} - \frac{\dd \nu}{\dd(\mu+\nu)}   \right| }\, \dd (\mu+\nu) = 2 \sup_{A\subseteq \X, \text{mbl}} |\mu(A)-\nu(A)|.
\end{align*}
Alternatively $TV$ can be written as a Wasserstein distance
\begin{align*}\label{eq:TVviaOT}
TV(\mu, \nu)= 2 \inf_{\pi\in \cpl(\mu, \nu)} \int \rho(x,y)\, \dd\pi(x,y),    
\end{align*}
where $\rho(x,y)=I_{x\neq y}$ denotes the discrete metric and $\cpl(\mu, \nu)$ the set of couplings or transport plans  i.e.\ probabilities on $\X\times \X$ with marginals $\mu, \nu$. 
Pinsker's classical inequality asserts that \[TV(\mu, \nu)\leq \sqrt{2H(\mu|\nu)},\] where $H$ denotes the relative entropy, i.e.\ $H(\mu|\nu) =\int\log(\dd \mu / \dd \nu)\, \dd\mu$ if $\mu\ll \nu$ and $H(\mu|\nu)= \infty $ otherwise. 

Our goal is to provide an \emph{adapted} version, which is of interest in the case where $\mu$ and $\nu$ are the laws of stochastic processes $(X_k)_{k=1}^n$ and $(Y_k)_{k=1}^n$, respectively. Here, $\X$ will be of product type, and it is notationally convenient to assume that $\X = \X_1 \times \ldots \times \X_n$, where each $\X_i$ is a Polish space in which $X_i$ and $Y_i$, for $i \leq n$, take values. When measuring the distance between stochastic processes, it is often useful to consider a variant of the Wasserstein distance that takes into account the inherent information structure of the processes. The resulting \emph{adapted Wasserstein distance} has a variety of applications from stochastic optimization, stochastic control, and mathematical finance to the theory of geometric inequalities and machine learning; see \cite{PfPi14, La18, XuAc21, BaBePa21, AcKraPa24, CoLi24, JiOb24, BaWi23} and the references therein.
The difference to the usual Wasserstein distance is that one restricts to transport plans which are adapted / causal in the following sense: Let $(\F_k)_{k\leq n}$ be the canonical filtration on $\X$, i.e.\ $\F_k$ is generated by the Borel sets of the form $A_1\times \ldots \times A_k\times \X_{k+1}\times \ldots \times \X_n\subseteq \X$ and consider the  $\mu$-disintegration $(\pi_x)_x$ of $\pi\in \cpl(\mu, \nu)$. Then 
\begin{align}\label{eq:CausalDef} \pi \text{ is \emph{causal} } :\Longleftrightarrow \ 
 x \mapsto \pi_x(B) \text{  is $\F_k$-measurable for each $B\in \F_k$.}
 \end{align}
 %I.e.\ the amount of mass transported from $x=(x_1, \ldots, x_n)$ to a set which depends only on the first $k$ coordinates is already determined by $x_1, \ldots, x_k$. 
When $\pi$ is of Monge-type, i.e.\ induced by some $T=(T_k)_{k=1}^n: \X\to \X$ it is causal if and only if $(T_k)_{k=1}^n$ is an adapted process, see \cite[Remark 2.3]{BaBeLiZa17}. 
A transport plan $\pi$ is \emph{bicausal}, in signs $\pi\in \cplbc(\mu, \nu)$, if \eqref{eq:CausalDef} holds for the $\mu$- and the $\nu$-disintegration.  %for the set of bicausal couplings. Adapted Wasserstein distances differ from their usual counterparts in that the minimization is restricted to $\cplbc(\mu, \nu)$. 
The \emph{adapted total variation} (see \cite{EcPa22}) is then given by 
\begin{align*}%\label{eq:ATV}
ATV(\mu, \nu):= 2 \inf_{\pi\in \cplbc(\mu, \nu)} \int \rho(x,y)\, \dd\pi(x,y) \quad \big(\ \geq TV(\mu, \nu) \ \big) . 
\end{align*}
Our main result is the following adapted Pinsker inequality: 
\begin{theorem}\label{thm:MainTheorem}
  Assume that $\mu, \nu$ are probabilities on the Polish space $\X_1\times \ldots \times \X_n$. Then
  \begin{align}\label{eq:AP} 
  ATV(\mu, \nu)\leq \sqrt{n}\sqrt{2H(\mu|\nu)}.\end{align}
\end{theorem}
As in the classical case, this inequality is tight, see \Cref{cor:PinskerTight} below. 

\pagebreak

\medskip

We provide some context on the adapted Pinsker inequality \eqref{eq:AP}. 

The adapted total variation was first studied by Eckstein--Pammer \cite{EcPa22}. They establish that  $ATV(\mu, \nu)\leq (2^n -1) TV(\mu, \nu)$,  %\begin{align}\label{eq:EcPa_bound}ATV(\mu, \nu)\leq (2^n -1) TV(\mu, \nu),\end{align}
in particular the two distances are equivalent. 
This is remarkable since as soon as one of the spaces $(\X_k, d_k), k < n$ is not discrete, the resulting Wasserstein distance $W$ and its adapted counterpart $AW$ are not equivalent, in fact $\AW$ generates strictly finer topologies. 
The idea of bounding an adapted distance by its classical counterpart $W$ under appropriate regularity conditions is put center stage in the work of Blanchet, Larsson, Park and Wiesel \cite{BlLaPaWi24}. In particular they establish that  if one restricts to the set of probabilities with whose kernels are sufficiently regular,  then $AW(\mu, \nu)\leq C^n \sqrt{W(\mu, \nu)}$.

%\begin{align}\label{eq:Wi_bound}AW(\mu, \nu)\leq C^n W(\mu, \nu).\end{align}
In view of the  dependence on the number of periods in these bounds,  the constant $\sqrt{n}$ in the adapted Pinsker inequality \eqref{eq:AP} appears  surprisingly small. 

On the other hand it can be expected that the chain rule  of entropy is well suited to control adapted variants of the Wasserstein distance. To give an example in this direction, recall that a probability $\gamma$ is said to satisfy a $T_2$-inequality with constant $C$ if the quadratic Wasserstein distance $W_2$ satisfies $W_2(\mu, \gamma) \leq \sqrt{2C H(\mu|\gamma)}$ for all probabilities $\mu$. From chain rule of entropy it follows easily that $T_2$ inequalities tensorize perfectly, i.e.\ $W_2(\mu, \otimes_{k=1}^n\gamma) \leq \sqrt{2C H(\mu|\otimes_{k=1}^n\gamma)}$, see e.g.\ \cite[Proposition 22.5]{Vi09}. Indeed this argument actually implies that the same inequality holds also for the adapted Wasserstein distance, i.e.\ $$AW_2(\mu, \otimes_{k=1}^n\gamma) \leq \sqrt{2C H(\mu|\otimes_{k=1}^n\gamma)},$$ see \cite[Section 5.2]{BaBeLiZa17}  for details. Strikingly $T_2$-inequalities exists also for adapted Wasserstein distances in continuous time (with the Wiener measure as reference measure), see the works of Föllmer \cite{Fo22a} and Lasalle \cite{La18}.




\section{Proofs}
As a preparation for the proof, we require an alternative characterization of bicausal couplings. It will be convenient to use shorthands like $x_{1:k}:=(x_1,\ldots,x_k) $.
Furthermore, we consider the successive  disintegrations
\begin{align*}
    \mu\left(\mathrm{d}x_{1:n}\right) &= \mu_1\left(\mathrm{d}x_1\right)  \mu^{x_1}\left(\mathrm{d}x_2 \right) \ldots \mu^{x_{1:(n-1)}} \left(\mathrm{d}x_n \right) \\
    \nu\left(\mathrm{d}y_{1:n}\right) &= \nu_1\left(\mathrm{d}y_1\right)  \nu^{y_1}\left(\mathrm{d}y_2 \right) \ldots \nu^{y_{1:(n-1)}} \left(\mathrm{d}y_n \right) \\
    \pi\left(\mathrm{d}(x_{1:n},y_{1:n})\right) &= \pi_1\left(\mathrm{d}(x_1,y_1)\right)  \pi^{x_1,y_1}\left(\mathrm{d}(x_2,y_2) \right) \ldots \pi^{x_{1:(n-1)},y_{1:(n-1)}} \left(\mathrm{d}(x_n,y_n) \right).
\end{align*}    
Then, as shown in \cite[Proposition 5.1]{BaBeLiZa17},  $\pi\in \cplbc(\mu, \nu)$ if and only if 
\begin{align}\label{eq:iteratechar}
   \pi_1\in \cpl(\mu_1, \nu_1) \quad \text{and} \quad \pi^{x_{1:k},y_{1:k} }\in \cpl(\mu^{x_{1:k}}, \nu^{y_{1:k}}), \text{ a.s. for $k< n$.} 
\end{align}
This alternative characterization is useful since it allows to calculate adapted transport problems in an iterated fashion.
Specifically, we use it to obtain the following representation of adapted total variation.\footnote{We note that this was observed independently by Acciaio, Hou, and Pammer (oral communication).}
\begin{lemma}\label{le:ATVsimple}
    %Suppose $X_1,\ldots, X_n$ are polish spaces and let $\mu, \nu \in \mathcal{P}(X)$, where $X:= X_1 \times \ldots \times X_n$. Then we have the following representation of the adapted total variation between measures $\mu,\nu \in \mathcal{P}(X)$
   The adapted total variation distance can be written as
    \begin{align*}
        ATV(\mu,\nu) & = TV(\mu_1,\nu_1) + \int_{\X_1} TV(\mu^{x_1}, \nu^{x_1}) \, \mathrm{d}(\mu_1 \wedge \nu_1)(x_1) + \ldots + \\
        & \int_{\X_{n-1}} TV(\mu^{x_{1:(n-1)}},\nu^{x_{1:(n-1)}}) \, \mathrm{d}(\mu^{x_{1:(n-2)}} \wedge \nu^{x_{1:(n-2)}})(x_{(n-1)}) \ldots \mathrm{d}(\mu_1 \wedge \nu_1)(x_1).
    \end{align*}
    Moreover the infimum in the definition of $ATV$ is attained.
    Indeed, any coupling $\pi$ for which $\pi_1$ and all $ \pi^{x_{1:k},y_{1:k}}, k< n$ put the maximal possible mass on the diagonal is optimal. 
%    holds.
\end{lemma}
Recall that minimum of two probabilities $\mu, \nu$ is given by
 \begin{align*}
\mu\wedge \nu= %\sup_{|f|\leq 1} \int f\, \dd (\mu-\nu) =
%|\mu\vee \nu|-|\mu\wedge\nu|.
\int{\textstyle  \left| \frac{\dd \mu}{\dd(\mu+\nu)} \wedge \frac{\dd \nu}{\dd(\mu+\nu)}   \right| }\, \dd (\mu+\nu) .
\end{align*}
\begin{proof}[Proof of \Cref{le:ATVsimple}]
    We interpret $\rho$ as the function $I_{x\neq y}$, irrespective of the space containing $x, y$ and note that $\rho(x_{1:k},y_{1:k}) = \rho(x_1,y_1) \vee \rho(x_2,y_2) \vee \ldots  \vee \rho(x_k,y_k)$. %In context it will always be clear on which space the function $\rho$ is defined. 
   
    
    For $n=1$, there is nothing to prove. For illustrative purposes we first show the result   in the case $n=2$.
    \begin{flalign*}
        & \phantom{=} ATV(\mu,\nu) \\
        & = 2 \inf_{\pi \in \cplbc(\mu,\nu)} \int_{\X \times \X} \!\!{\rho\left(x_{1:2},y_{1:2}\right)} \ \mathrm{d}\pi\left(x_{1:2},y_{1:2} \right) &&\\
        & = 2 \inf_{\pi_1 \in \cpl(\mu_1,\nu_1)} \ \inf_{\pi^{x_1,y_1} \in \cpl(\mu^{x_1},\nu^{y_1})} \int_{\X_1 \times \X_1} \int_{\X_2 \times \X_2} \!\! \rho(x_1,y_1) \vee \rho(x_2,y_2) \ \mathrm{d}\pi^{x_1,y_1}(x_2,y_2) \ \mathrm{d}\pi_1(x_1,y_1) &&\\ 
        & = 2 \inf_{\pi_1 \in \cpl(\mu_1,\nu_1)} \ \int_{\X_1 \times \X_1} \!\! \rho(x_1,y_1) \ \vee \ \inf_{\pi^{x_1,y_1} \in \cpl(\mu^{x_1},\nu^{y_1})} \int_{\X_2 \times \X_2} \!\! \rho(x_2,y_2) \ \mathrm{d}\pi^{x_1,y_1}(x_2,y_2) \ \mathrm{d}\pi_1(x_1,y_1) &&\\ % Fallunterscheidungen ob \rho(x_1,y_1) gleich 1 oder 0 ist
        & = \inf_{\pi_1 \in \cpl(\mu_1,\nu_1)} \ \int_{\X_1 \times \X_1} \!\! 2\rho(x_1,y_1) \ \vee \  TV(\mu^{x_1},\nu^{y_1}) \ \mathrm{d}\pi_1(x_1,y_1) &&\\ % Definition Totalvariation
        & = TV(\mu_1, \nu_1) + \int_{\X_1} \! TV(\mu^{x_1},\nu^{x_1}) \ \mathrm{d}(\mu_1 \wedge \nu_1)(x_1) % Den Schritt verstehe ich nicht 
    \end{flalign*}
To pass from line 2 to line 3,  we use the characterization of bicausal couplings given in \eqref{eq:iteratechar} together with a measurable selection arguments, see \cite[Propositon 5.2]{BaBeLiZa17} for a detailed justification.  To pass from line 5 to line 6, note that in line 5 we are given an ordinary transport problem with respect to a cost function that is $\leq 2$ on the diagonal and takes the value $2$ otherwise. A transport plan is optimal  for this problem if it puts as much mass on the diagonal as possible. From this we obtain the desired equality. Note also that the argument shows that all appearing infima are attained if the $\pi_1$ and the kernels $\pi^{x_1,y_1}$ put the maximal possible mass on the diagonal. In particular the infimum in the definition of $ATV$ is attained.
    Thus, we have proved the result in the case $n=2$. 
    
    The general case follows by induction. By  hypothesis we have
    \begin{align*}
        & 2 \inf_{\pi^{x_1,y_1} \in \cpl(\mu^{x_1},\nu^{y_1})} \int_{\X_2 \times \X_2} \!\! \rho(x_2,y_2)  \vee  \ldots  \vee \\ 
        &\quad \inf_{\pi^{x_{1:n-1},y_{1:n-1}} \in \cpl(\mu^{x_{1:n-1},y_{1:n-1}})} \int_{\X_n \times \X_n} \!\! \rho(x_n,y_n) \ \mathrm{d} \pi^{x_{1:n-1},y_{1:n-1}} (x_n,y_n) \ldots \mathrm{d}\pi^{x_1,y_1}(x_1,y_1) \notag \\
        & = TV(\mu^{x_1},\nu^{y_1}) + \int_{\X_2 \times \X_2} \!\! TV(\mu^{x_{1:2}}, \nu^{x_{1:2}}) \ + \ \ldots \ \\ 
        & + \int_{\X_{n-1}\times \X_{n-1}} \!\!\!\! TV(\mu^{x_{1:n-1}}, \nu^{x_{1:n-1}}) \ \mathrm{d}(\mu^{x_{1:n-2}} \wedge \nu^{x_{1:n-2}}) \ldots \mathrm{d}(\mu^{x_1} \wedge \nu^{x_1})(x_2). \notag
    \end{align*}
    Thus, it follows
    \begin{align*}
        & \phantom{=} ATV(\mu,\nu) \\
        & = 2 \inf_{\pi \in \cplbc(\mu,\nu)} \int_{\X \times \X} \!\! \rho(x_{1:n},y_{1:n}) \ \mathrm{d} \pi(x_{1:n},y_{1:n}) \\
        & = 2 \inf_{\pi_1 \in \cpl(\mu_1, \nu_1)} \int_{\X_1 \times \X_1} \rho(x_1,y_1) \vee \inf_{\pi^{x_1,y_1} \in \cpl(\mu^{x_1},\nu^{y_1})} \int_{\X_2 \times \X_2} \!\! \rho(x_2,y_2) \vee \ldots \vee \\ 
        & \inf_{\pi^{x_{1:n-1},y_{1:n-1}} \in \cpl(\mu^{x_{1:n-1},y_{1:n-1}})} \int_{\X_n \times \X_n} \rho(x_n,y_n) \ \mathrm{d} \pi^{x_{1:n-1},y_{1:n-1}} (x_n,y_n) \ldots \mathrm{d}\pi^{x_1,y_1}(x_1,y_1) \mathrm{d}\pi_1(x_1,y_1) \notag \\ % Iteriert wie im Fall n=2 aufteilen
        & = \inf_{\pi_1 \in \cpl(\mu_1, \nu_1)} \int_{\X_1 \times \X_1} \!\! 2\rho(x_1,y_1) \vee  \Big[ TV(\mu^{x_1},\nu^{y_1}) + \int_{\X_2 \times \X_2} \!\! TV(\mu^{x_{1:2}}, \nu^{x_{1:2}}) \ + \ \ldots \ \\ 
        & + \int_{\X_{n-1}\times \X_{n-1}} \!\!\!\! TV(\mu^{x_{1:n-1}}, \nu^{x_{1:n-1}}) \ \mathrm{d}(\mu^{x_{1:n-2}} \wedge \nu^{x_{1:n-2}}) \ldots \mathrm{d}(\mu^{x_1} \wedge \nu^{x_1})(x_2) \Big] \mathrm{d}   \pi_1(x_1,y_1) \notag \\
        & = TV(\mu_1,\nu_1) + \int_{\X_1} \!\! TV(\mu^{x_1}, \nu^{x_1}) \ + \ldots + \\
        & \int_{\X_{n-1}} \!\! TV(\mu^{x_{1:n-1}},\nu^{x_{1:n-1}}) \ \mathrm{d}(\mu^{x_{1:n-2}} \wedge \nu^{x_{1:n-2}})(x_{n-1}) \ldots \mathrm{d}(\mu_1 \wedge \nu_1)(x_1). \notag % Der Schritt ist mir ähnlich wie im Fall n=2 nicht klar
    \end{align*}
    Attainment follows as before.    This concludes the proof.
\end{proof}



\begin{proof}[Proof of \Cref{thm:MainTheorem}]
    
    For the case $n=1$, this is just the regular Pinsker inequality, and thus there is nothing to prove. For illustrative purposes, we will carry out the proof in the case $n=2$ before we prove the general bound.
Using  \Cref{le:ATVsimple} we find
    \begin{flalign*}
        & \phantom{=} ATV(\mu,\nu) \\
        & = TV(\mu_1, \nu_1) + \int_{\X_1} \!\! TV(\mu^{x_1},\nu^{x_1}) \ \mathrm{d}(\mu_1 \wedge \nu_1)(x_1) && \\ % Den Schritt verstehe ich nicht %MB: 
        & \leq \sqrt{2H(\mu_1|\nu_1)} + \int_{\X_1} \!\! \sqrt{2H(\mu^{x_1}|\nu^{x_1})} \mathrm{d}(\mu_1 \wedge \nu_1)(x_1) && \\
        & \leq \sqrt{1+|\mu_1 \wedge \nu_1|} \ \sqrt{2H(\mu_1|\nu_1) + \int_{\X_1} \!\! 2H(\mu^{x_1}|\nu^{x_1}) \mathrm{d}(\mu_1 \wedge \nu_1)(x_1)} &&\\
        & \leq \sqrt{2} \ \sqrt{2H(\mu_1|\nu_1) + \int_{\X_1} \!\! 2H(\mu^{x_1}|\nu^{x_1}) \mathrm{d}\mu_1(x_1)} &&\\
        & = \sqrt{2} \ \sqrt{2H(\mu|\nu)}.
    \end{flalign*}
    The first inequality follows from Pinsker's inequality, the second follows from defining the measure $m:= \delta_{0} + \mu_1 \wedge \nu_1$ and the function 
    $$g(s):= \begin{cases}
        2 H(\mu_1|\nu_1) & \text{if }   s=0 \\
        2 H(\mu^{x_1}|\nu^{x_1}) & \text{if} \ s \in \X_1
    \end{cases}$$ 
    and then applying Jensen's inequality to obtain $$\left(\int_{\{0\} \bigcup \X_1} \!\! \sqrt{g} \ \mathrm{d}m / \| m\|  \right)^2\leq \int_{\{0\} \bigcup \X_1} \!\! g \ \mathrm{d}m / \| m\|.$$ 
    The last equality is established by applying the chain rule for entropy.
    The case of an arbitrary $n$ will be analogous to the case $n=2$:
    \begin{align*}
        & \phantom{=} ATV(\mu,\nu) \\ 
        & = TV(\mu_1,\nu_1) + \int_{\X_1} \! TV(\mu^{x_1}, \nu^{x_1}) + \ldots + \int_{\X_{n-1}} \!\!\!\!\! TV(\mu^{x_{1:n-1}},\nu^{x_{1:n-1}}) \ \mathrm{d}(\mu^{x_{1:n-2}} \wedge \nu^{x_{1:n-2}}) \ldots \mathrm{d} (\mu_1 \wedge \nu_1) \\ 
        & \leq \sqrt{n} \ \sqrt{2 H(\mu_1|\nu_1) + \ldots + \int_{\X_1 \times \ldots \times \X_{n-1}} \!\!\!\! 2 H(\mu^{x_{1:n-1}},\nu^{x_{1:n-1}}) \mathrm{d}(\mu^{x_{1:n-2}} \wedge \nu^{x_{1:n-2}}) \ldots \mathrm{d} (\mu_1 \wedge \nu_1}) \\
        & \leq \sqrt{n} \ \sqrt{2 H(\mu|\nu)}.
    \end{align*}
    The first inequality follows from first applying the Pinsker inequality and then defining a measure $m:= \delta_{t} + \mu_1 \wedge \nu_1 + (\mu_1 \wedge \nu_1) \otimes (\mu^{x_1} \wedge \nu^{x_1}) \ + \ldots + (\mu_1 \wedge \nu_1) \otimes (\mu^{x_1} \wedge \nu^{x_1}) \otimes \ldots \otimes (\mu^{x_{1:n-1}} \wedge \nu^{x_{1:n-1}})$ and a function $g$ similarly as in the case $n=2$ on $\{0\} \bigcup \X_1 \bigcup \ldots \bigcup \X_1 \times \ldots \times \X_{n-1}$ and applying the Jensen inequality. Note that since each measure in the definition of $m$ has a total measure less than $1$, we have $|m| \leq n$. The last inequality follows from applying the chain rule of the entropy for disintegration repeatedly.
\end{proof}

The following lemma is of course well known, but we give the proof here since we will use the specific construction in the proof to establish the tightness of the adapted Pinsker inequality \eqref{eq:AP}.
\begin{lemma}\label{tightpinsker}
    The Pinsker inequality is tight.
\end{lemma}

\begin{proof}
    Consider a measurable space $(\Omega = \{0,1\},2^\Omega\}$ and measures $\nu = \frac{1}{2}\delta_0+\frac{1}{2}\delta_1$ and $\mu = (\frac{1}{2}+\varepsilon)\ \delta_0 + (\frac{1}{2}-\varepsilon)\ \delta_1$ for $\varepsilon \in \left(0,\frac{1}{2}\right)$. It follows that $TV(\mu,\nu)=2\varepsilon$. 
    Furthermore, using the Taylor expansion of the logarithm, we obtain 
    \begin{align*}\textstyle H(\mu|\nu) & = \textstyle (\frac{1}{2}+\varepsilon) \log\left(\frac{(\frac{1}{2}+\varepsilon)}{\frac{1}{2}}\right)+(\frac{1}{2}-\varepsilon)\log\left(\frac{(\frac{1}{2}-\varepsilon)}{\frac{1}{2}}\right)\\
    &= \textstyle (\frac{1}{2}+\varepsilon) \log(1+2\varepsilon)+(\frac{1}{2}-\varepsilon)\log(1-2\varepsilon) \\ & =  \textstyle (\frac{1}{2}+\varepsilon)(2\varepsilon-2\varepsilon^2)+(\frac{1}{2}-\varepsilon)(-2\varepsilon-2\varepsilon^2)+o(\varepsilon^2) \textstyle = 2\varepsilon^2 + o(\varepsilon^2).
    \end{align*}
    Thus, we conclude $\lim\limits_{\varepsilon \searrow 0} \frac{TV(\mu,\nu)^2}{H(\mu,\nu)}=\lim\limits_{\varepsilon \searrow 0} \frac{(2\varepsilon)^2}{2\varepsilon^2+o(\varepsilon^2)} =2$, proving the tightness of the Pinsker inequality.
\end{proof}

\begin{corollary}\label{cor:PinskerTight}
    The adapted Pinsker inequality $ATV(\mu,\nu) \leq \sqrt{n} \sqrt{2 H(\mu|\nu)}$ is tight for every $n \in \mathbb{N}$. % Ich nutze \mathbb{N}={1,2,3,...}
\end{corollary}
    
\begin{proof}
    The case $n=1$ is exactly \Cref{tightpinsker}, thus we assume that $n \ge 2$ for the rest of the proof. Consider the measurable space $(\Omega:=\{0,1\}^n,2^\Omega)$ and measures $\mu_i=(\frac{1}{2}+\varepsilon)\ \delta_0 + (\frac{1}{2}-\varepsilon)\ \delta_1$ and $\nu_i = \frac{1}{2}\delta_0+\frac{1}{2}\delta_1$ on $\{0,1\}$ for every $i\in \{1,\ldots,n\}$ as in \Cref{tightpinsker}. Further, define the product measures $\mu := \bigotimes_{i=1}^n \mu_i$ and $\nu := \bigotimes_{i=1}^n \nu_i$. Note that the disintegration of $\mu$ is given by $\mu=\underbrace{\mu_1 \otimes \cdots \otimes \mu_1}_\text{$n$-times}$ 
    and analogously for $\nu$.

    In order to determine $ATV(\mu,\nu)$ we use the decomposition from \Cref{le:ATVsimple} and the fact, that $|\mu_1 \wedge \nu_1| = 1-\varepsilon$ to obtain
    \begin{align*}
        & \phantom{=} ATV(\mu,\nu) \\
        & = TV(\mu_1,\nu_1)+\int_{\X_1} TV(\mu_1,\nu_1)+\ldots +\int_{\X_{n-1}} TV(\mu_1,\nu_1) \ \mathrm{d}(\mu_1 \wedge \nu_1)(x_{n-1}) \ldots \mathrm{d}(\mu_1 \wedge \nu_1)(x_1) \\
        & = 2\varepsilon \sum_{k=0}^{n-1}(1-\varepsilon)^k  = 2\varepsilon \frac{1-(1-\varepsilon)^n}{\varepsilon}  = 2-2(1-\varepsilon)^n.
    \end{align*}
    Since $\mu$ and $\nu$ are product measures, it follows similarly as in the proof of \Cref{tightpinsker} that $H(\mu,|\nu) = n H(\mu_1|\nu_1) = n (2\varepsilon^2+o(\varepsilon^2))$. Due to \Cref{thm:MainTheorem} we  already know that $\frac{ATV(\mu,\nu)^2}{nH(\mu|\nu)}\leq2$ or equivalently $\frac{(2-2(1-\varepsilon)^n)^2}{n^2(2\varepsilon^2+o(\varepsilon^2))} \leq 2$. We conclude the proof by showing that the left hand side of the inequality converges to $2$ as $\varepsilon$ approaches $0$:
    \begin{align*}
        & \phantom{=} \lim_{\varepsilon \searrow 0} \frac{(2-2(1-\varepsilon)^n)^2}{2n^2\varepsilon^2+o(\varepsilon^2)} \\
        & = \lim_{\varepsilon \searrow 0} \frac{4-8(1-\varepsilon)^n+4(1-\varepsilon)^{2n}}{2n^2\varepsilon^2+o(\varepsilon^2)} \\ 
        & = \lim_{\varepsilon \searrow 0} \frac{4-8+8n\varepsilon-8\frac{n(n-1)}{2}\varepsilon^2+o(\varepsilon^2)+4-8n\varepsilon+4\frac{2n(2n-1)}{2}\varepsilon^2+o(\varepsilon^2)}{2n^2\varepsilon^2+o(\varepsilon^2)} \\
        & = \lim_{\varepsilon \searrow 0} \frac{4n^2\varepsilon^2+o(\varepsilon^2)}{2n^2\varepsilon^2+o(\varepsilon^2)} 
         = \lim_{\varepsilon \searrow 0} \frac{4n^2+\frac{o(\varepsilon^2)}{\varepsilon^2}}{2n^2+\frac{o(\varepsilon^2)}{\varepsilon^2}} 
         = \frac{4n^2}{2n^2} = 2.
    \end{align*}
    Note that the same argument applies provided that each space $\X_k$ contains a set of measure $1/2$.
\end{proof}



\medskip
\noindent
{\bf Acknowledgment:}
 This research was funded  in whole or in part by the Austrian Science Fund (FWF) [doi: 10.55776/P35197]. For open access purposes, the author has applied a CC BY public copyright license to any author-accepted manuscript version arising from this submission.	
    
\bibliographystyle{abbrv} 
\bibliography{joint_biblio}

\end{document}

